\chapter{Methodology}\label{chap:met}

This chapter describes the research methodology, system design, and architectural decisions used to evaluate Kubernetes as a platform for implementing business logic. The chapter expands upon the high-level methodology outlined in Chapter~\ref{chap:intro}, Section~\ref{sec:methodology}, providing a detailed description of both the Kubernetes-based system and the traditional baseline system.

\section{Research Approach}

This research followed a system design and comparative analysis approach. Two functionally equivalent systems were designed and implemented:

\begin{itemize}
    \item A Kubernetes-native system based on Custom Resource Definitions (CRDs) and the Operator pattern.
    \item A traditional backend system based on a REST API and relational database.
\end{itemize}

Both systems implemented the same business domain logic to ensure comparability. The primary difference lies in how the business logic is modeled, executed, and managed.

The methodology focuses on analyzing:
\begin{itemize}
    \item Architectural differences between declarative and imperative approaches.
    \item Execution models (event-driven reconciliation vs request-response).
    \item Observability and system introspection capabilities.
\end{itemize}

\section{Business Domain Modeling}

The selected domain for this research is a billing and subscription management system. This domain was chosen because it includes both stateful entities and time-based business logic, making it suitable for evaluating different architectural paradigms.

\subsection{Core Entities}

The system defines two primary entities:

\begin{itemize}
    \item \textbf{BillingPlan}: Represents a pricing tier, including attributes such as price, currency, billing period, and usage limits.
    \item \textbf{Subscription}: Represents a user's subscription to a billing plan, including references to the plan, lifecycle state, and billing timestamps.
\end{itemize}

\subsection{Business Logic}

The implemented business logic includes:

\begin{itemize}
    \item \textbf{Activation}: Upon creation of a subscription, the system validates the referenced billing plan, activates the subscription, and initializes billing timestamps.
    \item \textbf{Billing Cycle Management}: The system tracks the last successful payment and the scheduled next billing time, advancing both whenever a billing cycle completes.
    \item \textbf{Termination}: A subscription leaves the active state either because the user explicitly ends it or because the system detects a terminal condition (a missed billing deadline or a payment failure).
\end{itemize}

The two systems implement the same conceptual lifecycle but expose it through different mechanisms, and this asymmetry is itself a finding of the comparison. Table~\ref{tab:state-comparison} enumerates the observable states in each implementation.

\begin{table}[h]
\centering
\small
\caption{Subscription lifecycle states in each implementation}
\label{tab:state-comparison}
\begin{tabularx}{\textwidth}{@{}lXX@{}}
\toprule
\textbf{Stage} & \textbf{backend-billing} & \textbf{kube-billing} \\ \midrule
Initial & \texttt{Active} (set on creation) & \texttt{Active} (set by first reconciliation) \\
User-initiated end & \texttt{Canceled} (via \texttt{POST /subscriptions/\{id\}/cancel}) & resource deletion guarded by a finalizer; a pro-rated refund event is emitted before removal \\
Missed billing deadline & \texttt{Expired} (set by background worker once per minute) & not modelled as a state; the reconciler keeps re-attempting billing \\
Payment processing error & not modelled & \texttt{PaymentError} (price parse or charge failure) \\
Referenced plan missing & request rejected at admission (404) & \texttt{Error} with condition \texttt{BillingPlanNotFound = True} \\ \bottomrule
\end{tabularx}
\end{table}

The asymmetry reflects an architectural truth: in a REST-and-database design, cancellation and expiration are in-band state transitions that must be modelled explicitly; in a controller-based design, cancellation can ride on Kubernetes' built-in resource lifecycle via finalizers, while transient failures are surfaced through Conditions and re-attempted by the reconcile loop rather than being persisted as terminal states. Chapter~\ref{chap:eval} discusses how this distinction affects the comparative evaluation.

\section{Traditional System Architecture}

The baseline system follows a conventional monolithic architecture implemented in Go.

\subsection{Technology Stack}

\begin{itemize}
    \item \textbf{Programming language}: Go 1.25
    \item \textbf{HTTP framework}: Gorilla Mux (lightweight router with middleware support)
    \item \textbf{Database}: PostgreSQL 15 (relational database with ACID guarantees)
    \item \textbf{ORM}: GORM v2 (Go ORM with automatic migration support)
    \item \textbf{Configuration}: Viper (environment variables + YAML files)
    \item \textbf{Validation}: go-playground/validator v10
    \item \textbf{Observability}: OpenTelemetry SDK + Prometheus client\_golang
    \item \textbf{Rate Limiting}: Token bucket algorithm (10 req/s, burst 20)
\end{itemize}

\subsection{Architecture Overview}

The system implements a layered architecture with clear separation of concerns:

\begin{itemize}
    \item \textbf{Handlers layer}: HTTP request processing and response formatting
    \item \textbf{Repository layer}: Database operations through GORM interfaces
    \item \textbf{Models layer}: GORM entities (BillingPlan, Subscription)
    \item \textbf{Workers layer}: Background subscription expiration checker
    \item \textbf{Middleware layer}: Rate limiting, panic recovery, tracing
\end{itemize}

The REST API exposes endpoints for billing plans and subscriptions. Business logic is implemented within HTTP handlers and repository methods, with explicit error handling and validation.

\subsection{Execution Model}

The system uses a synchronous request-response model with asynchronous background processing:

\begin{itemize}
    \item \textbf{Synchronous path}: Clients send HTTP requests $\rightarrow$ handlers validate input $\rightarrow$ repository executes database operations $\rightarrow$ response returned
    \item \textbf{Asynchronous path}: Background worker runs every 30 seconds, queries PostgreSQL for expired subscriptions, updates state to \texttt{Expired}
\end{itemize}

The background worker implements time-based business logic that cannot be handled synchronously:

\begin{lstlisting}[language=Go, caption={Background expiration checker worker}, label={lst:worker}]
// Worker runs every 30 seconds
ticker := time.NewTicker(30 * time.Second)
for range ticker.C {
    expiredSubs := repo.FindExpired(time.Now())
    for _, sub := range expiredSubs {
        sub.State = "Expired"
        repo.Update(sub)
    }
}
\end{lstlisting}

\subsection{Data Persistence}

All system state is stored in PostgreSQL with the following schema:

\begin{itemize}
    \item \textbf{BillingPlan table}: id (PK), name (unique), price (string for precision), currency, billing\_period, created\_at, updated\_at, deleted\_at (soft delete)
    \item \textbf{Subscription table}: id (PK), user\_id, plan\_ref (FK to BillingPlan.name), state (Active/Cancelled/Expired), last\_payment, next\_billing, created\_at, updated\_at, deleted\_at
\end{itemize}

GORM auto-migration ensures schema synchronization on application startup. The database serves as the single source of truth, with business logic tightly coupled to database operations through the repository pattern.

\subsection{API Endpoints}

\begin{itemize}
    \item \texttt{GET /plans} --- List all billing plans
    \item \texttt{POST /plans} --- Create billing plan (validates name, price, currency)
    \item \texttt{GET /subscriptions} --- List all subscriptions
    \item \texttt{GET /subscriptions/\{id\}} --- Get subscription by ID
    \item \texttt{POST /subscriptions} --- Create subscription (auto-activates)
    \item \texttt{POST /subscriptions/\{id\}/cancel} --- Cancel active subscription
    \item \texttt{GET /health} --- Health check endpoint
    \item \texttt{GET /ready} --- Readiness check (includes database connectivity)
    \item \texttt{GET /metrics} --- Prometheus metrics endpoint
\end{itemize}

\section{Kubernetes-Based System Architecture}

The experimental system is implemented using Kubernetes as a platform for business logic through Custom Resource Definitions and the Operator pattern.

\subsection{Technology Stack}

\begin{itemize}
    \item \textbf{Programming language}: Go 1.25.3
    \item \textbf{Operator Framework}: Kubebuilder v4.13.0 + controller-runtime v0.23.1
    \item \textbf{CRD API Group}: \texttt{billing.cloud-native.io/v1alpha1}
    \item \textbf{Metrics}: Prometheus client\_golang v1.23.1
    \item \textbf{Tracing}: OpenTelemetry SDK (OTLP over HTTP, port 4318)
    \item \textbf{Testing}: Ginkgo v2 + Gomega + envtest (local API server)
    \item \textbf{E2E Testing}: Kind (Kubernetes in Docker)
\end{itemize}

\subsection{Custom Resource Definitions}

Two CRDs were defined to extend the Kubernetes API:

\begin{itemize}
    \item \texttt{BillingPlan} (\texttt{billingplans.billing.cloud-native.io}) --- defines pricing configuration
    \item \texttt{Subscription} (\texttt{subscriptions.billing.cloud-native.io}) --- represents user subscription
\end{itemize}

\subsubsection{BillingPlan Specification}

The \texttt{BillingPlan} CRD includes comprehensive validation:

\begin{itemize}
    \item \texttt{spec.price} (string, required): Pattern \texttt{\textasciicircum{}\textbackslash d+(\textbackslash.\textbackslash d\{1,2\})?\$}, min length 1, max length 20 --- ensures decimal with up to 2 decimal places
    \item \texttt{spec.currency} (string, required): Enum (USD, EUR, RUB, KZT, GBP, JPY, CNY)
    \item \texttt{spec.billingPeriod} (string, required): Enum (hourly, daily, weekly, monthly, yearly)
    \item \texttt{spec.requeueIntervalSeconds} (integer, optional): Range 1--31536000 (1 year), default 30 --- controls reconciliation frequency
    \item \texttt{spec.limits} (map, optional): Key-value pairs for usage constraints (e.g., \texttt{apiCalls: 10000})
\end{itemize}

\textbf{Status subresource} (managed by controller):
\begin{itemize}
    \item \texttt{status.activeSubscriptions} (integer): Count of active subscriptions referencing this plan
    \item \texttt{status.totalRevenue} (string): Cumulative revenue from all subscriptions
    \item \texttt{status.conditions} (array): Standard Kubernetes conditions (\texttt{Available}, etc.)
\end{itemize}

\subsubsection{Subscription Specification}

The \texttt{Subscription} CRD enforces strict validation:

\begin{itemize}
    \item \texttt{spec.userId} (string, required): Pattern \texttt{\textasciicircum{}[a-zA-Z0-9\_-]+\$}, max length 256 --- alphanumeric with underscores and hyphens
    \item \texttt{spec.planRef} (string, required): Pattern \texttt{\textasciicircum{}[a-z0-9]([-a-z0-9]*[a-z0-9])?\$}, max length 253 --- Kubernetes naming convention
\end{itemize}

\textbf{Status subresource} (managed by controller):
\begin{itemize}
    \item \texttt{status.state} (string): The current lifecycle state. The controller writes one of \texttt{Active}, \texttt{PaymentError}, or \texttt{Error}; the CRD schema does not constrain the field with an enum so that future state values can be introduced without a CRD version bump.
    \item \texttt{status.lastPayment} (timestamp): Time of last successful billing.
    \item \texttt{status.nextBilling} (timestamp): Scheduled time for next billing cycle.
    \item \texttt{status.observedGeneration} (integer): Tracks status freshness against \texttt{spec}.
    \item \texttt{status.conditions} (array of \texttt{metav1.Condition}): Standardised condition types \texttt{Active}, \texttt{BillingPlanNotFound}, \texttt{PaymentError}, each with \texttt{True/False/Unknown} status, a machine-readable \texttt{reason}, and a human-readable \texttt{message}.
\end{itemize}

\subsection{Operator Controller Implementation}

The core business logic resides in the \texttt{SubscriptionReconciler}, implemented using controller-runtime. The reconciler continuously processes Subscription resources and ensures correct system state.

\subsubsection{Reconciliation Workflow}

The reconciliation logic follows a deterministic 6-step sequence:

\begin{enumerate}
    \item \textbf{Fetch Subscription}: Retrieve the Subscription resource from the API server
    \item \textbf{Handle Deletion}: If \texttt{DeletionTimestamp.IsZero() == false}:
    \begin{itemize}
        \item Execute final billing (pro-rated refund calculation --- TODO in current implementation)
        \item Remove finalizer \texttt{billing.cloud-native.io/finalizer}
        \item Return, allowing Kubernetes to complete deletion
    \end{itemize}
    \item \textbf{Add Finalizer}: If finalizer absent, attach it to prevent premature deletion
    \item \textbf{Validate BillingPlan}: Query for referenced BillingPlan by \texttt{planRef}:
    \begin{itemize}
        \item If not found: Set \texttt{State = Error}, \texttt{BillingPlanNotFound = True}, requeue after 60 seconds
        \item Emit Kubernetes Event: \texttt{Warning BillingPlanNotFound}
    \end{itemize}
    \item \textbf{Initialize Subscription}: If \texttt{State == ""} (new subscription):
    \begin{itemize}
        \item Set \texttt{State = Active}
        \item Increment Prometheus metric \texttt{billing\_active\_subscriptions}
        \item Set \texttt{LastPayment = now}, \texttt{NextBilling = now + RequeueIntervalSeconds}
        \item Set conditions: \texttt{Active = True}, \texttt{BillingPlanNotFound = False}, \texttt{PaymentError = False}
        \item Emit Event: \texttt{Normal SubscriptionActivated}
    \end{itemize}
    \item \textbf{Process Billing Cycle}: If \texttt{State == Active} and \texttt{now > NextBilling}:
    \begin{itemize}
        \item Increment Prometheus metric \texttt{billing\_revenue\_total} by plan price
        \item Update \texttt{LastPayment = now}, \texttt{NextBilling = now + RequeueIntervalSeconds}
        \item Emit Event: \texttt{Normal PaymentProcessed}
        \item \textbf{RequeueAfter}: \texttt{RequeueIntervalSeconds} from BillingPlan.spec
    \end{itemize}
\end{enumerate}

\subsubsection{Watches and Event Sources}

The controller watches multiple event sources:

\begin{itemize}
    \item \textbf{Primary}: \texttt{Subscription} resources (create, update, delete events)
    \item \textbf{Secondary}: \texttt{BillingPlan} resources via \texttt{mapPlanToSubscriptions} function --- triggers reconciliation of all Subscriptions referencing a changed Plan
\end{itemize}

This watch pattern enables automatic propagation of plan changes (e.g., price updates) to all affected subscriptions.

\subsubsection{Finalizer Handling}

A Kubernetes finalizer (\texttt{billing.cloud-native.io/finalizer}) ensures safe deletion:

\begin{itemize}
    \item When user executes \texttt{kubectl delete subscription sub-user1}, Kubernetes sets \texttt{DeletionTimestamp} but does not remove the resource
    \item Reconciler detects deletion timestamp, executes final billing logic
    \item After finalizer removal, Kubernetes completes resource deletion
\end{itemize}

This mechanism guarantees cleanup logic execution even if the controller crashes during processing.

\subsection{Execution Model}

Unlike the traditional system, execution is event-driven and declarative:

\begin{itemize}
    \item \textbf{Declarative}: User declares desired state in YAML (\texttt{spec}), controller drives actual state (\texttt{status})
    \item \textbf{Event-driven}: Resource changes trigger reconciliation events via Kubernetes watch API
    \item \textbf{Asynchronous}: Controllers process events concurrently with configurable worker count (default: 1)
    \item \textbf{Eventually consistent}: System converges toward desired state through repeated reconciliation
\end{itemize}

No direct client-driven execution of business logic occurs --- the controller autonomously manages state transitions.

\subsection{Data Persistence}

System state is stored in the Kubernetes control plane (\texttt{etcd}) as custom resources:

\begin{itemize}
    \item \textbf{etcd}: Distributed key-value store (backing Kubernetes API server)
    \item \textbf{CRD instances}: YAML resources persisted with full versioning and metadata
    \item \textbf{Status subresource}: Separate update path for controller-managed state (optimizes write conflicts)
\end{itemize}

The declarative state model enables automatic recovery: if the controller crashes, it reconstructs state from etcd on restart and resumes reconciliation.

\subsection{Observability Integration}

\subsubsection{Prometheus Metrics}

Three custom business metrics are exposed at \texttt{:8080/metrics}:

\begin{itemize}
    \item \texttt{billing\_active\_subscriptions} (Gauge): Current count of Active subscriptions
    \item \texttt{billing\_revenue\_total} (Counter): Cumulative revenue processed (sum of all successful billing cycles)
    \item \texttt{billing\_payment\_failures} (Counter): Number of failed payment attempts (defined but not incremented in current implementation)
\end{itemize}

Additionally, controller-runtime exposes standard metrics:
\begin{itemize}
    \item \texttt{controller\_runtime\_reconcile\_total} (Counter): Total reconciliation attempts by result (success, error, requeue)
    \item \texttt{controller\_runtime\_reconcile\_errors\_total} (Counter): Total reconciliation errors
\end{itemize}

\subsubsection{OpenTelemetry Tracing}

Distributed tracing captures asynchronous reconciliation flow:

\begin{itemize}
    \item \textbf{Root span}: \texttt{ReconcileSubscription} --- covers entire reconciliation cycle
    \item \textbf{Child span}: \texttt{ProcessPayment} --- billing logic execution
    \item \textbf{Mapped span}: \texttt{MapPlanToSubscriptions} --- triggered by BillingPlan watch events
\end{itemize}

Traces are exported via OTLP over HTTP to \texttt{localhost:4318} (Jaeger-compatible).

\subsubsection{Kubernetes Events}

The controller emits audit events for significant state changes:

\begin{itemize}
    \item \texttt{Normal SubscriptionActivated}: Subscription successfully activated
    \item \texttt{Normal PaymentProcessed}: Billing cycle completed
    \item \texttt{Warning PaymentFailed}: Payment processing failed
    \item \texttt{Normal FinalBilling}: Final billing executed on deletion
\end{itemize}

Events are accessible via \texttt{kubectl get events --field-selector involvedObject.name=sub-user1}.

\subsection{Deployment Configuration}

The system is deployed using production-ready Kubernetes manifests:

\begin{itemize}
    \item \textbf{HorizontalPodAutoscaler (HPA)}: Min 1, Max 5 replicas, CPU target 80\%, Memory target 80\%
    \item \textbf{PodDisruptionBudget (PDB)}: Min available 1 --- ensures availability during voluntary disruptions
    \item \textbf{Resource limits}: CPU 100m--500m, Memory 128Mi--256Mi
    \item \textbf{NetworkPolicy}: Restricts \texttt{/metrics} endpoint access to Prometheus namespace
    \item \textbf{RBAC}: ServiceAccount with minimal permissions (get, list, watch Subscriptions and BillingPlans)
\end{itemize}

Deployment is managed via Kustomize overlays (\texttt{config/default}) combining CRDs, RBAC, controller Deployment, and monitoring configuration.

\section{Comparison of Architectural Approaches}

The two systems differ fundamentally in how they implement business logic. Table~\ref{tab:arch-comparison} summarizes the key architectural differences.

\begin{table}[h]
\centering
\small
\caption{Architectural Comparison: Traditional Backend vs Kubernetes Operator}
\label{tab:arch-comparison}
\begin{tabularx}{\textwidth}{@{}lXX@{}}
\toprule
\textbf{Aspect} & \textbf{backend-billing} & \textbf{kube-billing} \\ \midrule
State storage & PostgreSQL (external) & etcd via CRD (control plane) \\
API style & REST HTTP (port 8080) & Kubernetes API (kubectl) \\
Execution model & Request-response + background worker & Reconcile loop + RequeueAfter \\
Scaling & Manual (Docker Compose) & HPA (automatic, 1--5 replicas) \\
Fault tolerance & Depends on DB + orchestrator & Kubernetes self-healing \\
Validation & Runtime (go-playground/validator) & Admission time (CRD schema) \\
Error handling & Explicit (error returns) & Eventual consistency + retries \\
Observability & HTTP metrics + DB tracing & Business metrics + Events + Conditions \\
Deployment & Docker container & Kubernetes Deployment + HPA + PDB \\
Configuration & Environment variables + YAML & ConfigMap + Secrets \\ \bottomrule
\end{tabularx}
\end{table}

\subsection{State Management}

\textbf{backend-billing} stores state in PostgreSQL, an external relational database. This approach provides:
\begin{itemize}
    \item Strong ACID guarantees for transactions
    \item Mature query capabilities (JOINs, aggregations, indexes)
    \item Separation of application and data layers
\end{itemize}

\textbf{kube-billing} stores state in etcd via Custom Resources. This approach provides:
\begin{itemize}
    \item Declarative state management (spec vs status separation)
    \item Built-in versioning and metadata (generation, resourceVersion)
    \item Automatic persistence through Kubernetes control plane
\end{itemize}

\subsection{Execution Flow}

\textbf{backend-billing} follows imperative execution:
\begin{lstlisting}[language=bash, caption={Traditional REST API request flow}, label={lst:rest-flow}]
POST /subscriptions → validate → INSERT → 201 Created
\end{lstlisting}

\textbf{kube-billing} follows declarative reconciliation:
\begin{lstlisting}[language=bash, caption={Kubernetes Operator reconciliation flow}, label={lst:k8s-flow}]
kubectl apply -f sub.yaml → etcd write → watch event → Reconcile() → status update
\end{lstlisting}

\subsection{Time-Based Logic}

\textbf{backend-billing} uses a ticker-based worker:
\begin{itemize}
    \item Fixed interval (30 seconds)
    \item Polls database for expired subscriptions
    \item Updates state synchronously within worker goroutine
\end{itemize}

\textbf{kube-billing} uses RequeueAfter:
\begin{itemize}
    \item Dynamic interval from BillingPlan.spec.requeueIntervalSeconds
    \item Controller-runtime schedules next reconciliation
    \item No polling --- event-driven wake-up
\end{itemize}

\subsection{Lifecycle Management}

\textbf{backend-billing} deletes immediately:
\begin{itemize}
    \item \texttt{DELETE /subscriptions/\{id\}} → soft delete (deleted\_at timestamp)
    \item No cleanup hook execution guarantee
\end{itemize}

\textbf{kube-billing} uses finalizers:
\begin{itemize}
    \item \texttt{kubectl delete} → DeletionTimestamp set → finalizer blocks deletion
    \item Reconciler executes final billing → removes finalizer → Kubernetes deletes
    \item Guaranteed cleanup even on controller failure
\end{itemize}

\section{Observability Design}

Observability was implemented in both systems to enable comparative analysis of system behavior, with emphasis on business-level events rather than only infrastructure metrics.

\subsection{Kubernetes System Observability}

The Kubernetes-based system integrates three observability layers:

\subsubsection{Prometheus Metrics}

Custom business metrics exposed at \texttt{:8080/metrics}:
\begin{itemize}
    \item \texttt{billing\_active\_subscriptions} (Gauge) --- current count of Active subscriptions
    \item \texttt{billing\_revenue\_total} (Counter) --- cumulative revenue from successful billing cycles
    \item \texttt{billing\_payment\_failures} (Counter) --- failed payment attempts (defined, not incremented)
\end{itemize}

Standard controller-runtime metrics:
\begin{itemize}
    \item \texttt{controller\_runtime\_reconcile\_total} --- reconciliation attempts by result
    \item \texttt{controller\_runtime\_reconcile\_errors\_total} --- error count
    \item \texttt{workqueue\_depth} --- pending work items
    \item \texttt{process\_cpu\_seconds}, \texttt{go\_memstats\_alloc\_bytes} --- runtime metrics
\end{itemize}

\subsubsection{OpenTelemetry Tracing}

Span hierarchy for reconciliation:
\begin{itemize}
    \item \texttt{ReconcileSubscription} (root) --- entire reconciliation cycle
    \item \texttt{ProcessPayment} (child) --- billing logic execution
    \item \texttt{MapPlanToSubscriptions} (child) --- triggered by BillingPlan watch
\end{itemize}

Export: OTLP over HTTP to \texttt{localhost:4318} (Jaeger-compatible).

\subsubsection{Kubernetes Events}

Audit trail via Events API:
\begin{itemize}
    \item \texttt{Normal SubscriptionActivated} --- new subscription activated
    \item \texttt{Normal PaymentProcessed} --- billing cycle completed
    \item \texttt{Warning PaymentFailed} --- payment error
    \item \texttt{Normal FinalBilling} --- final billing on deletion
\end{itemize}

\subsection{Traditional System Observability}

The traditional system implements two-layer observability:

\subsubsection{Prometheus Metrics}

HTTP and business metrics at \texttt{/metrics}:
\begin{itemize}
    \item \texttt{billing\_subscriptions\_created\_total} (Counter)
    \item \texttt{billing\_subscriptions\_cancelled\_total} (Counter)
    \item \texttt{http\_request\_duration\_seconds} (Histogram) --- request latency
    \item \texttt{http\_requests\_total} (Counter) --- request count by method/path
\end{itemize}

\subsubsection{OpenTelemetry Tracing}

Integrated tracing for HTTP and database layers:
\begin{itemize}
    \item HTTP middleware (\texttt{otelmux}) --- traces request lifecycle
    \item GORM plugin --- traces SQL queries (duration, statement, rows affected)
\end{itemize}

Export: OTLP over HTTP to Jaeger (\texttt{http://jaeger:4318}).

\subsection{Comparative Analysis}

The Kubernetes system provides richer business-level observability:
\begin{itemize}
    \item \textbf{Events} --- audit trail without custom logging
    \item \textbf{Conditions} --- standardized state reporting in resource status
    \item \textbf{Reconciliation metrics} --- visibility into control loop behavior
\end{itemize}

The traditional system provides simpler but more direct observability:
\begin{itemize}
    \item \textbf{HTTP metrics} --- immediate request/response visibility
    \item \textbf{Database tracing} --- direct query-level insights
\end{itemize}

\section{Experimental Methodology}

Experimental evaluation was conducted in a controlled local environment to compare the two architectural approaches across multiple dimensions.

\subsection{Evaluation Environment}

Both systems were deployed and tested on identical hardware. Full reproducibility details, including workload generator parameters and the number of repeated trials, are documented in Section~\ref{sec:reproducibility}.

\begin{itemize}
    \item \textbf{CPU}: Apple M3 (8 cores: 4 performance + 4 efficiency)
    \item \textbf{RAM}: 8\,GB unified memory
    \item \textbf{OS}: macOS 15.7.7 Sequoia (Darwin 24.6.0)
    \item \textbf{Kubernetes}: Minikube v1.33.1 (for kube-billing)
    \item \textbf{Docker}: Docker Desktop v27.5.0 (for backend-billing)
\end{itemize}

\subsection{Evaluation Metrics}

The evaluation focuses on four categories:

\subsubsection{Code Metrics (Lean Software Development)}

\begin{itemize}
    \item \textbf{Lines of Code (LOC)} --- total Go source lines (excluding tests)
    \item \textbf{File count} --- number of Go source files
    \item \textbf{Cyclomatic complexity} --- maximum complexity per function
    \item \textbf{Test coverage} --- percentage of code covered by automated tests
    \item \textbf{Dependencies} --- number of Go module dependencies
\end{itemize}

\subsubsection{Performance Metrics}

\begin{itemize}
    \item \textbf{Latency} --- p50, p95, p99 response times for CRUD operations
    \item \textbf{Throughput} --- requests per second under load
    \item \textbf{Resource usage} --- CPU (millicores) and memory (MB) consumption
    \item \textbf{Startup time} --- time from launch to readiness
\end{itemize}

\subsubsection{Operational Metrics}

\begin{itemize}
    \item \textbf{Scaling} --- manual vs automatic (HPA)
    \item \textbf{Fault tolerance} --- recovery from controller/database failure
    \item \textbf{Deployment complexity} --- number of configuration files, steps
\end{itemize}

\subsubsection{Observability Coverage}

\begin{itemize}
    \item \textbf{Metrics count} --- number of custom business metrics
    \item \textbf{Trace coverage} --- percentage of code paths traced
    \item \textbf{Event/audit trail} --- presence of audit logging
\end{itemize}

\subsection{Benchmark Methodology}

\subsubsection{Load Testing}

\textbf{backend-billing}:
\begin{lstlisting}[language=bash, caption={Load testing commands for backend-billing}, label={lst:load-backend}]
hey -n 1000 -c 10 http://localhost:8080/plans
hey -n 1000 -c 50 http://localhost:8080/subscriptions
\end{lstlisting}

\textbf{kube-billing}:
\begin{lstlisting}[language=bash, caption={Load testing commands for kube-billing}, label={lst:load-kube}]
for i in {1..1000}; do kubectl apply -f plan.yaml; done
kubectl apply -f subscription.yaml (sequential, rate-limited)
\end{lstlisting}

\subsubsection{Resource Monitoring}

\begin{itemize}
    \item \textbf{backend-billing}: \texttt{docker stats}, Go \texttt{pprof}
    \item \textbf{kube-billing}: \texttt{kubectl top pods}, Prometheus metrics
\end{itemize}

\subsubsection{Microbenchmarks}

Go testing package for isolated function benchmarks:
\begin{lstlisting}[language=bash, caption={Go microbenchmark commands}, label={lst:bench}]
go test -bench=BenchmarkCreatePlan -benchmem
go test -bench=BenchmarkValidatePlan -benchmem
\end{lstlisting}

\subsection{Threats to Validity}

Several limitations affect the generalizability of results:

\begin{itemize}
    \item \textbf{Local environment} --- results may differ in production Kubernetes clusters with network latency and resource contention
    \item \textbf{Single-node testing} --- no distributed system effects (network partitions, leader election)
    \item \textbf{Simulated payments} --- no external payment gateway integration (latency, failures)
    \item \textbf{Short test duration} --- no long-term etcd compaction or database bloat effects
\end{itemize}

Despite these limitations, the comparative analysis provides meaningful insights into architectural trade-offs.

\section{Summary}

This chapter described the methodology used to evaluate Kubernetes as a platform for business logic. Two systems implementing identical functionality were designed using different paradigms: a traditional REST-based system with PostgreSQL and a Kubernetes-native Operator-based system with CRDs. The architectural differences --- imperative vs declarative, request-response vs reconciliation, external database vs etcd --- provide the foundation for the comparative evaluation presented in Chapter~\ref{chap:eval}.
